\chapter{Poda con NLF}
\label{cap:prune-nlf}

\begin{intuicion}
NLF (\emph{Neighborhood Label Frequency}) es el primer prune del
matcher: una tabla \([\V \times \Lc_V] \to \mathbb{N}\) que cuenta,
para cada vértice y cada tipo destino, cuántos vecinos salientes
tiene de ese tipo. Antes de iniciar el DFS desde un \emph{start}
\(s\), el matcher chequea: ¿el primer paso de la hipótesis pide
``vecino de tipo \(L\)''? Si \(\NLF(s, L) = 0\), \(s\) no puede
generar ningún match. Lo descarta. Es local (\(O(1)\) por start),
sound (nunca descarta un start válido), y barato de mantener
(capítulo~\ref{cap:incremental-nlf}). Este capítulo formaliza el
prune, prueba su soundness, y muestra cuánto poda en grafos
típicos.
\end{intuicion}

\section{Definición formal}

\begin{definition}[NLF]
Sea \(G = (\V, \E)\) un grafo dirigido temporal etiquetado, y
\(L \in \Lc_V\) un \texttt{EntityType} con \emph{tag} en
\([0, \texttt{NLF\_TYPE\_COUNT})\). Para \(v \in \V\):
\[
\NLF(v, L) = |\{w \in N^{+}(v) : \mathrm{type}(w) = L\}|.
\]
\end{definition}

\noindent
Es decir: \(\NLF(v, L)\) cuenta los vecinos salientes de \(v\) cuyo
tipo es exactamente \(L\). Es una matriz \(|\V| \times
|\Lc_V|\) de enteros (en GraphHunter, \texttt{u32}).

\begin{nota}
\(\NLF\) ignora el \emph{tipo de arista} y el \emph{timestamp}.
Sólo cuenta vecinos por tipo destino. Eso lo hace
poco-discriminante para hipótesis con muchos predicados, pero
\emph{universalmente útil} para descartar starts donde \emph{no
existe el tipo destino}.
\end{nota}

\section{El prune}

\begin{algorithm}[H]
\caption{Prune NLF al inicio del DFS}
\KwIn{starts candidatos \(\V_{L_0^V}\), hipótesis \(H\)}
\KwOut{starts filtrados}
filtered \(\leftarrow \emptyset\)\;
\For{\(s \in \V_{L_0^V}\)}{
    \If{\texttt{satisfies\_first\_step}(s, H)}{
        filtered.push(s)\;
    }
}
\Return filtered\;
\BlankLine
\Comment{\texttt{satisfies\_first\_step(s, H)}}
\(L \leftarrow L_1^V\) (tipo destino del primer paso de \(H\))\;
\Return \(\NLF(s, L) \geq 1\)\;
\end{algorithm}

\noindent
La versión más completa: para hipótesis con \emph{shape} graph,
donde un mismo vértice puede tener varios pasos salientes
(estrella centrada en él), el chequeo se generaliza a
\(\NLF(s, L_j^V) \geq r_j\) para cada hijo \(j\) requerido.

\section{Soundness}

\begin{theorem}[NLF prune sound]
Si \(\NLF(s, L) = 0\), entonces \(s\) no es origen de ningún match
de la hipótesis \(H\) cuyo primer paso requiere un vecino de tipo
\(L\).
\end{theorem}

\begin{proof}
Por contradicción. Supongamos que existe un match
\(\sigma\) con \(\sigma(p_0) = s\). Por definición de match, existe
\(w = \sigma(p_1) \in \V\) con \(\mathrm{type}(w) = L_1^V = L\) y
una arista \((s, w, \tau) \in \E\). Es decir,
\(w \in N^{+}(s)\) y \(\mathrm{type}(w) = L\). Entonces
\(\NLF(s, L) \geq 1\), contradicción. \(\square\)
\end{proof}

\noindent
\textbf{Importante}: la otra dirección NO vale. \(\NLF(s, L) \geq
1\) NO garantiza que \(s\) genere match. El prune es
\emph{conservador}: nunca descarta starts válidos, pero puede
\emph{aceptar} starts que después no generen matches (falsos
positivos). El DFS los procesa, no encuentra nada, retorna sin
emit.

\subsection{Prune más fuerte: chequear varios tipos}

Si la hipótesis tiene shape graph, el primer vértice puede tener
\emph{varios} hijos (capítulo~\ref{cap:dsl-graph}):

\begin{lstlisting}
User c -[Auth]-> Host h1; c -[Auth]-> Host h2; c -[Auth]-> Host h3
\end{lstlisting}

\noindent
Acá el start \(c\) necesita \emph{tres} hijos de tipo \texttt{Host}.
NLF se generaliza: \(\NLF(c, \texttt{Host}) \geq 3\). El prune
descarta cualquier \(c\) con menos vecinos del tipo correcto.

\subsection{NLF en pasos intermedios}

¿Y por qué no chequear NLF en cada paso del DFS, no sólo al
inicio? Porque el DFS \emph{ya} itera \(N^{+}(v)\) en cada paso ---el
chequeo NLF sería redundante salvo en el primer paso, donde nos
ahorra entrar al loop. En los pasos intermedios el costo de iterar
y filtrar por tipo es comparable al costo del chequeo NLF.

\section{Cuánto poda en la práctica}

Tres escenarios reales:

\subsection{Escenario 1: hipótesis con primer paso muy selectivo}

Hipótesis: \texttt{User -[Auth]-> Host}. Grafo: \(10^5\) usuarios,
de los cuales sólo \(10^4\) tienen aristas \texttt{Auth}
salientes hacia \texttt{Host}s.

\begin{itemize}
    \item Sin NLF: el DFS arranca desde \(10^5\) starts; los
    \(9 \cdot 10^4\) que no tienen aristas Host generan iteración
    vacía sobre \(N^{+}(s)\) y retornan.
    \item Con NLF: descarta los \(9 \cdot 10^4\) antes de entrar al
    DFS. \(10\times\) menos starts a procesar.
\end{itemize}

\subsection{Escenario 2: starts ``hojas''}

Grafos scale-free tienen muchos vértices de grado bajo. Para
hipótesis con primer paso de tipo \texttt{Host}, los \emph{User}s
de grado 0 (no se loguearon en ninguna sesión analizada) están en
\(\V_{User}\) pero NLF los descarta inmediatamente.

\begin{itemize}
    \item Típico: \(\sim 30\%\) de usuarios sin aristas en una
    sesión semanal (gente de vacaciones, cuentas inactivas, etc.).
    NLF los filtra a costo cero.
\end{itemize}

\subsection{Escenario 3: hipótesis multi-hijo (shape graph)}

Hipótesis: \texttt{c -[Auth]-> Host h1; c -[Auth]-> Host h2; c -[Auth]-> Host h3}
(usuario que se autentica en al menos 3 hosts distintos).

\begin{itemize}
    \item Sin NLF generalizada: el DFS arranca desde todos los
    \texttt{User}; descubre on-the-fly que la mayoría sólo tiene 1
    o 2 hosts.
    \item Con NLF generalizada: \(\NLF(u, \texttt{Host}) \geq 3\).
    En grafos típicos, sólo el \(\sim 5\%\) de usuarios cumple. El
    DFS arranca desde \(20\times\) menos starts.
\end{itemize}

\section{Implementación: \texttt{NlfTable}}

El detalle de implementación está en el capítulo~\ref{cap:nlf-detalle}
(layout, \texttt{IndexBacking}, mmap-spill). Acá lo que importa
para el prune es la firma:

\begin{lstlisting}[language=Rust]
impl NlfTable {
    pub fn count(&self, sid: StrId, type_tag: u8) -> u32 {
        let row = &self.counts.get(sid.index());  // [u32; NLF_TYPE_COUNT]
        row[type_tag as usize]
    }

    pub fn can_match(
        &self,
        sid: StrId,
        required: &[(u8, u32)],  // (type_tag, min_count)
    ) -> bool {
        required.iter().all(|&(tag, min)| self.count(sid, tag) >= min)
    }
}
\end{lstlisting}

\noindent
\texttt{can\_match} es la operación que el matcher invoca al inicio
del DFS para cada start candidato. \(O(|\text{required}|)\) que
es \(\leq |\Lc_V| = 9\). Despreciable.

\begin{enelcodigo}
\filepath{core/graph-engine/src/nlf.rs}: \texttt{NlfTable::can\_match}.
El \emph{backing} es \texttt{IndexBacking<[u32; 9]>} ---heap-resident
para grafos chicos, \emph{mmap-spilled} arriba del \emph{budget}
(default 1\,GB).
\end{enelcodigo}

\section{Construcción}

\begin{algorithm}[H]
\caption{NLF build}
\KwIn{grafo \(G = (\V, \E)\)}
\KwOut{\texttt{NlfTable}}
\(\mathrm{counts} \leftarrow \mathrm{Vec\,de\,}|\V|\,\mathrm{filas}\,[u32; \texttt{NLF\_TYPE\_COUNT}],\) inicializado a 0\;
\For{\((u, v, \tau) \in \E\)}{
    \(L \leftarrow \mathrm{type}(v).\mathrm{tag}\)\;
    \If{\(L < \texttt{NLF\_TYPE\_COUNT}\)}{
        \(\mathrm{counts}[u][L] \leftarrow \mathrm{counts}[u][L] + 1\)\;
    }
}
\Return \texttt{NlfTable\{ counts \}}\;
\end{algorithm}

\noindent
\(O(|\E|)\) iteraciones, una pasada secuencial por
\texttt{edge\_store}. Para \(|\E| = 10^8\), \(\sim 30\) segundos en
hardware típico ---el costo que el \emph{lazy build} difiere al
primer hunt (capítulo~\ref{cap:finalize}).

\section{Limitaciones}

NLF es un prune simple y por eso tiene limitaciones:

\begin{itemize}
    \item \textbf{No considera tipo de arista.} Si la hipótesis
    pide \texttt{User -[Auth]-> Host} y un \texttt{User} tiene
    aristas \texttt{Spawn} pero no \texttt{Auth} a algún
    \texttt{Host}, NLF no lo descarta (\(\NLF(u, \texttt{Host}) \geq
    1\) por el Spawn). El DFS verifica el tipo de arista en cada
    paso e ignora.
    \item \textbf{No considera ventana temporal.} Si todas las
    aristas \texttt{Auth} de \(u\) están fuera del \texttt{WITHIN},
    NLF no lo descarta. Para esto sirve el k-hop reachability
    extendido temporal (capítulo~\ref{cap:prune-khop}).
    \item \textbf{Solo el primer paso.} El prune actúa antes de
    iniciar el DFS, no en los pasos intermedios. Para los pasos
    intermedios, el DFS itera \(N^{+}\) y filtra por tipo en el
    loop ---no hay forma de pre-filtrar sin gastar memoria.
\end{itemize}

\noindent
Cada limitación se ataca con una técnica adicional:
\emph{neighborhood label profile} (NLP, no implementado en
GraphHunter por costo) considera tipos de arista; k-hop temporal
considera ventana; \emph{label-and-edge histograms} considera tipos
intermedios.

\section{Resumen}

\begin{itemize}
    \item NLF: \(\NLF(v, L) = |\{w \in N^{+}(v) : \mathrm{type}(w)
    = L\}|\). Tabla \(|\V| \times |\Lc_V|\) de \texttt{u32}.
    \item Prune: descartar starts con \(\NLF(s, L) < r\). Sound
    (nunca descarta válidos), no completo (acepta falsos
    positivos).
    \item Build \(O(|\E|)\) en una pasada; cabe lazy en
    \texttt{OnceLock}.
    \item Limitaciones: ignora tipo de arista y ventana temporal.
    Sirve sólo para el primer paso.
    \item Impacto típico: 10\(\times\)--20\(\times\) menos starts a
    procesar; mayor en hipótesis multi-hijo.
\end{itemize}

\begin{ejercicios}
\begin{enumerate}
    \item \starred{} Para el grafo de la
    figura~\ref{fig:gt-ejemplo} (capítulo~\ref{cap:grafos-temporales}),
    construí la \(\NLF\) completa: una matriz con filas por vértice
    y columnas por \texttt{EntityType}.

    \item Para la hipótesis
    \texttt{User -[Auth]-> Host -[Connect]-> IP}, listá todos los
    starts (vértices \texttt{User}) que sobreviven al prune NLF en
    el grafo del ejercicio anterior.

    \item \starred{} Demostrá: si la hipótesis tiene shape graph
    con \emph{root} de tres hijos del mismo tipo \(L\), el prune
    correcto es \(\NLF(s, L) \geq 3\). ¿Y si los tres hijos son de
    tipos distintos?

    \item ¿Por qué NLF NO se chequea en pasos intermedios del DFS?
    Discutí: ¿qué información agregaría? ¿Cuál es el costo de
    chequearla?

    \item Diseñá un grafo y una hipótesis donde NLF tenga cero
    impacto (descarta cero starts). Después diseñá uno donde NLF
    descarte el 99\%.

    \item Si extendiéramos NLF para considerar también el tipo de
    arista (NLP: \emph{Neighborhood Label Profile}), ¿cuánta más
    memoria costaría? ¿Vale la pena?
\end{enumerate}
\end{ejercicios}
